\subsection{POSITIVE ROOTS}

Let C be a chamber of R, and let \(\mathrm{B}(\mathrm{C})=\{\alpha_{1},\ldots,\alpha_{l}\}\) be the corresponding basis of R. The order relation on V (resp. \(V^{*}\) ) defined by C is the order relation compatible with the vector space structure of V (resp. \(V^{*}\) ) for which the elements \(\geqslant0\) are the linear combinations of the \(\alpha_{i}\) (resp. the \(\alpha_{i}^{-}\) ) with coefficients \(\geqslant0\) . An element that is positive for one of these relations is said to be positive for C, or positive for the basis \(\mathrm{B}(\mathrm{C})\) . These order relations are also defined by \(C^{-}\) , as one sees by identifying V with \(V^{*}\) by using a scalar product invariant under \(\mathrm{W}(\mathrm{R})\) . In view of Th. 2, no. 5, an element of \(V^{*}\) is \(\geqslant0\) if and only if its values on C are \(\geqslant0\) . An element x of V is \(\geqslant0\) if and only if its values on \(C^{-}\) are \(\geqslant0\) , or, equivalently, if \((x|y)\geqslant0\) for all \(y\in C\) .

The elements of \(\overline{C}\) are \(\geqslant 0\) for \(C\) by Lemma 6 of Chap. V, § 3, no. 5. But the set of elements \(\geqslant 0\) for \(C\) is in general distinct from \(\overline{C}\) (cf.~Plate X, Systems \(A_2\) , \(B_2\) , \(G_2\) ).

THEOREM 3. Every root is a linear combination with integer coefficients of the same sign of elements of B(C). In particular, every root is either positive or negative for C.

If \(\alpha \in \mathbb{R}\) , the kernel \(L_{\alpha}\) of \(\alpha\) does not meet \(C^{\cdot}\) , so \(\alpha\) is either \(>0\) on the whole of \(C^{\cdot}\) or \(< 0\) on the whole of \(C^{\cdot}\) , hence the second assertion. It remains to show that \(\alpha\) is contained in the subgroup \(P\) of \(V\) generated by \(B(C)\) ; we can assume that \(\alpha\) is indivisible. Now the group P is clearly stable under the \(s_{\gamma}\) , for \(\gamma \in \mathrm{B}(\mathrm{C})\) , hence also under \(\mathrm{W}(\mathrm{R})\) by Th. 2. Since \(\alpha\) is of the form \(w(\beta)\) , with \(w \in \mathrm{W}(\mathrm{R})\) and \(\beta \in \mathrm{B}(\mathrm{C})\) (cf.~Prop. 15), we have \(\alpha \in P\) . Q.E.D.

Denote by \(R_{+}(C)\) the set of roots that are positive for C. Thus,

\[
R = R _ {+} (C) \cup (- R _ {+} (C))
\]

is a partition of \(R\) .

COROLLARY. Let \(\gamma\) be a linear combination of roots with integer coefficients, and \(\alpha\) an indivisible root. If \(\gamma\) is proportional to \(\alpha\) , then \(\gamma \in Z\alpha\) .

By Prop. 15 of no. 5, C can be chosen so that \(\alpha \in B(C)\) . By Th. 3,

\[
\gamma = \sum_ {\beta \in B (C)} n _ {\beta} \beta \text {with} n _ {\beta} \in \mathbf {Z}.
\]

Thus, if \(\gamma\) is proportional to \(\alpha\) , then \(\gamma = n_{\alpha}\alpha\) , which proves the corollary.

Now let S be the set of reflections \(s_{\alpha}\) for \(\alpha \in \mathrm{B}(\mathrm{C})\) and let T be the union of the conjugates of S under W. For \(\alpha \in \mathrm{B}(\mathrm{C})\) and \(w \in W\) , the element \(t = ws_{\alpha}w^{-1}\) of T is the orthogonal reflection \(s_{\beta}\) associated to the root \(\beta = w(\alpha)\) ; conversely, for any indivisible root \(\beta\) , there exists an element \(w \in W\) such that \(\alpha = w^{-1}(\beta) \in \mathrm{B}(\mathrm{C})\) (Prop. 15) and \(s_{\beta} = ws_{\alpha}w^{-1} \in T\) . It follows that a bijection \(\psi\) from the set of indivisible roots to \(\{\pm1\} \times T\) is obtained by associating to an indivisible root \(\beta\) the pair \((\varepsilon, s_{\beta})\) , where \(\varepsilon = +1\) if \(\beta\) is positive and \(\varepsilon = -1\) if \(\beta\) is negative.

On the other hand, (W, S) is a Coxeter system (Th. 2) and the results of Chap. IV, § 1, no. 4 can be applied. We have seen that, if w is an element of W of length (with respect to S) equal to q, there exists a subset \(T_{w}\) of T, with q elements, such that, if \(w = s_{1} \ldots s_{q}\) with \(s_{i} \in S\) and if

\[
t _ {i} = s _ {1} \dots s _ {i - 1} s _ {i} s _ {i - 1} \dots s _ {1}
\]

(for \(1 \leqslant i \leqslant q\) ), then \(\mathrm{T}_w = \{t_1, \ldots, t_q\}\) . Recall that we have also defined in no. 4 of § 1 a number \(\eta(w, t)\) (for \(w \in \mathrm{W}\) and \(t \in \mathrm{T}\) ) equal to \(+1\) if \(t \notin \mathrm{T}_w\) and to \(-1\) if \(t \in \mathrm{T}_w\) . Finally, recall that, if we define a map \(\mathrm{U}_w\) from the set \(\{\pm 1\} \times \mathrm{T}\) to itself by the formula

\[
\mathbf {U} _ {w} (\varepsilon , t) = (\varepsilon \eta (w ^ {- 1}, t), w t w ^ {- 1}),
\]

the map \(w \mapsto \mathrm{U}_w\) is a homomorphism from \(W\) to the group of permutations of the set \(\{\pm 1\} \times T\) (Chap. IV, §1, no. 4, Lemma 1).

PROPOSITION 17. Assume that R is reduced and let \(w \in W\) and \(\alpha \in R\) .

\begin{enumerate}
\def\labelenumi{(\roman{enumi})}
\item
  We have \(\psi(w(\alpha)) = U_w(\psi(\alpha)).\)
\item
  Assume that \(\alpha\) is positive. The root \(w(\alpha)\) is negative if and only if
\end{enumerate}

\[
\eta (w ^ {- 1}, s _ {\alpha}) = - 1,
\]

in other words if \(s_{\alpha} \in \mathrm{T}_{w^{-1}}\) .

\begin{enumerate}
\def\labelenumi{(\roman{enumi})}
\setcounter{enumi}{2}
\tightlist
\item
  We have \(\eta(w,s_{\alpha})=-1\) if and only if the chambers C and \(w(C)\) are on opposite sides of the hyperplane \(L_{\alpha}\) . In other words, the set \(T_{w}\) consists of the reflections with respect to the walls separating C and \(w(C)\) .
\end{enumerate}

Let \(\beta \in \mathrm{B}(\mathrm{C})\) and put \(s = s_{\beta}\) . Clearly \(\mathrm{T}_s = \{s\}\) and consequently

\[
\mathrm{U} _ {s} (\varepsilon , t) = \left\{ \begin{array}{l l} (\varepsilon , s t s ^ {- 1}) & \text { if } t \neq s \\ (- \varepsilon , s) & \text { if } t = s. \end{array} \right.\tag{12}
\]

On the other hand, let \(\rho = \sum_{\gamma \in B(C)} n_{\gamma}(\rho) \gamma\) be a positive root. Put

\[
s (\rho) = \sum_ {\gamma \in B (C)} n _ {\gamma} (s (\rho)) \gamma .
\]

If \(\rho \neq \beta\) , there exists an element \(\gamma \in \mathrm{B}(\mathbf{C})\) with \(\gamma \neq \beta\) , such that \(n_{\gamma}(\rho) > 0\) , and we have \(n_{\gamma}(s(\rho)) = n_{\gamma}(\rho) > 0\) (no. 1, formula (5)). Hence \(s(\rho)\) is positive. We deduce immediately that

\[
\psi (s (\varepsilon . \rho)) = \left\{ \begin{array}{l l} (\varepsilon , s s _ {\rho} s ^ {- 1}) & \text {if} \rho \neq \beta \\ (- \varepsilon , s) & \text {if} \rho = \beta . \end{array} \right.\tag{13}
\]

Comparison of (12) and (13) now shows that \(\mathrm{U}_{s}(\psi(\gamma)) = \psi(s(\gamma))\) for all roots \(\gamma\) and all \(s \in S\) . Since S generates W, (i) follows.

On the other hand, saying that \(w(\alpha)\) is negative is equivalent to saying that

\[
\psi (w (\alpha)) = (- 1, w s _ {\alpha} w ^ {- 1}),
\]

or, by (i), that \(\mathrm{U}_w(\psi (\alpha)) = (-1, ws_{\alpha}w^{-1})\) . If in addition \(\alpha\) is positive, then \(\psi (\alpha) = (+1,s_{\alpha})\) and \(\mathrm{U}_w(\psi (\alpha)) = (\eta (w^{-1},s_{\alpha}),ws_{\alpha}w^{-1})\) , hence (ii).

Finally, by (ii), \(\eta(w,s_{\alpha})=-1\) if and only if one of the roots \(\alpha\) and \(w^{-1}(\alpha)\) is positive and the other negative. This is equivalent to saying that \((\alpha|x).(w^{-1}(\alpha)|x)=(\alpha|x).( \alpha|w(x))<0\) for all \(x\in C\) , hence the first assertion in (iii). The second assertion in (iii) follows immediately.

COROLLARY 1. Let \(\beta \in \mathrm{B}(\mathbb{C})\) . The reflection \(s_{\beta}\) permutes the positive roots not proportional to \(\beta\) .

We reduce immediately to the case in which R is reduced. In that case, our assertion follows from (ii) and the fact that \(T_{s_{\beta}} = s_{\beta}\) .

COROLLARY 2. Assume that R is reduced. Let \(w \in W\) , let \(q\) be the length of \(w\) with respect to S (Chap. IV, § 1, no. 1), and let \(w = s_1 \ldots s_q\) be a reduced decomposition of \(w\) . Let \(\alpha_1, \ldots, \alpha_q\) be the elements of B(C) corresponding to \(s_1, \ldots, s_q\) . Put

\[
\theta_ {i} = s _ {q} s _ {q - 1} \dots s _ {i + 1} (\alpha_ {i}), \quad i = 1, \dots , q.
\]

The roots \(\theta_{i}\) are \(>0\) , pairwise distinct, \(w(\theta_{i}) < 0\) , and every root \(\alpha > 0\) such that \(w(\alpha) < 0\) is equal to one of the \(\theta_{i}\) .

Let \(X\) be the set of \(\alpha > 0\) such that \(w(\alpha) < 0\) . By (ii),

\[
\operatorname{Card} (X) = \operatorname{Card} \left(T _ {w ^ {- 1}}\right) = l \left(w ^ {- 1}\right) = l (w) = q.
\]

On the other hand, if \(\alpha \in X\) it is clear that there exists \(i \in [1, q]\) such that

\[
s _ {i + 1} \ldots s _ {q} (\alpha) > 0 \quad \text {and} \quad s _ {i} s _ {i + 1} \ldots s _ {q} (\alpha) <   0.
\]

By Cor. 1, this implies that \(s_i s_{i+1} \ldots s_q(\alpha) = \alpha_i\) and hence that \(\alpha = \theta_i\) . The set X is thus contained in the set of \(\theta_i\) . Since \(\text{Card}(X) = q\) , this is possible only if X is equal to the set of \(\theta_i\) and these are pairwise distinct. Hence the corollary.

COROLLARY 3. Assume that R is reduced. There exists a unique longest element \(w_0\) in W. Its length is equal to the number of positive roots and \(w_0\) transforms the chamber C to -C. We have \(w_0^2 = 1\) and \(l(ww_0) = l(w_0) - l(w)\) for all \(w \in W\) .

It is clear that -C is a chamber. Hence there exists an element \(w_{0}\) of W that transforms C to -C. Then \(w_{0}(\alpha) < 0\) for all positive roots \(\alpha\) and the first two assertions of Cor. 3 are immediate consequences of Cor. 2. We have \(w_{0}^{2}(C) = C\) , so \(w_{0}^{2} = 1\) . Finally, if \(w \in W\) , the length \(l(w)\) (resp. \(l(ww_{0})\) ) is equal, by Prop. 17 (iii), to the number of walls separating C and \(w(C)\) (resp. \(ww_{0}(C) = -w(C)\) ). Since \(w(C)\) and \(-w(C)\) are on opposite sides of every wall, the sum \(l(w) + l(w_{0})\) is equal to the total number of walls, that is to \(l(w_{0})\) .

PROPOSITION 18. Let \(x \in V\) . The following three properties are equivalent:

\begin{enumerate}
\def\labelenumi{(\roman{enumi})}
\item
  \(x \in \overline{C};\)
\item
  \(x \geqslant s_{\alpha}(x)\) for all \(\alpha \in \mathrm{B}(\mathrm{C})\) (with respect to the order relation defined by \(\mathrm{C}\) );
\item
  \(x \geqslant w(x)\) for all \(w \in W\) .
\end{enumerate}

Since \(s_{\alpha}(x) = x - \langle x, \alpha' \rangle \alpha\) and since \(\overline{\mathbf{C}}\) is the set of elements \(x \in V\) such that \(\langle x, \alpha' \rangle \geqslant 0\) for all \(\alpha \in \mathrm{B}(\mathbf{C})\) , the equivalence of (i) and (ii) is obvious. On the other hand, it is clear that (iii) \(\Longrightarrow\) (ii). We show that (i) \(\Longrightarrow\) (iii). Let \(x \in \overline{\mathbf{C}}\) , and let \(w \in W\) . We argue by induction on the length \(l(w)\) of \(w\) . The case \(l(w) = 0\) is trivial. If \(l(w) \geqslant 1\) , \(w\) can be written in the form \(w = w's_{\alpha}\) , with \(\alpha \in \mathrm{B}(\mathbf{C})\) and \(l(w') = l(w) - 1\) . Then

\[
x - w (x) = x - w ^ {\prime} (x) + w ^ {\prime} \left(x - s _ {\alpha} (x)\right).
\]

The induction hypothesis shows that \(x - w'(x)\) is positive. On the other hand,

\[
w ^ {\prime} (x - s _ {\alpha} (x)) = w \left(s _ {\alpha} (x) - x\right) = - \langle x, \alpha^ {*} \rangle w (\alpha).
\]

Now \(s_{\alpha} \in T_{w^{-1}}\) , and Prop. 17 (ii) shows that \(w(\alpha) < 0\) . Hence the result.

COROLLARY. An element \(x \in \mathbf{C}\) if and only if \(x > w(x)\) for all \(w \in \mathbf{W}\) such that \(w \neq 1\) .

PROPOSITION 19. Let \((\beta_{i})_{1\leqslant i\leqslant n}\) be a sequence of positive roots for the chamber C such that \(\beta_{1}+\beta_{2}+\cdots+\beta_{n}\) is a root. Then there exists a permutation \(\pi\in S_{n}\) such that, for all \(i\in\{1,2,\ldots,n\}\) , \(\beta_{\pi(1)}+\beta_{\pi(2)}+\cdots+\beta_{\pi(i)}\) is a root.

We argue by induction on n, the proposition being clear for \(n \leqslant 2\) . Put \(\beta = \beta_{1} + \cdots + \beta_{n}\) . Then \(\sum_{i=1}^{n} (\beta | \beta_{i}) = (\beta | \beta) > 0\) , so there exists an index k such that \((\beta | \beta_{k}) > 0\) . If \(\beta = \beta_{k}\) , then n = 1 since \(\beta_{i} > 0\) for all i. Otherwise \(\beta - \beta_{k}\) is a root (no. 3, Cor. of Th. 1); it then suffices to apply the induction hypothesis to \(\beta - \beta_{k} = \sum_{i \neq k} \beta_{i}\) .

COROLLARY 1. Let \(\alpha \in \mathrm{R}_{+}(\mathrm{C})\) . Then \(\alpha \in \mathrm{B}(\mathrm{C})\) if and only if \(\alpha\) is the sum of two positive roots.

If \(\alpha\) is the sum of two positive roots, Th. 3 shows that \(\alpha \in \mathrm{B}(\mathrm{C})\) . If \(\alpha \notin \mathrm{B}(\mathrm{C})\) , Th. 3 shows that \(\alpha = \sum_{k=1}^{n} \beta_k\) with \(\beta_k \in \mathrm{B}(\mathrm{C})\) for all \(k\) and \(n \geqslant 2\) . Permuting the \(\beta_k\) if necessary, we can assume that \(\sum_{k=1}^{n-1}\) is a root (Prop. 19), and hence that \(\alpha\) is the sum of the positive roots \(\sum_{k=1}^{n-1} \beta_k\) and \(\beta_n\) .

COROLLARY 2. Let \(\varphi\) be a map from \(R\) to a abelian group \(\Gamma\) having the following properties:

\begin{enumerate}
\def\labelenumi{\arabic{enumi})}
\item
  \(\varphi(-\alpha) = -\varphi(\alpha)\) for \(\alpha \in \mathbf{R}\) ;
\item
  if \(\alpha \in \mathrm{R}\) , \(\beta \in \mathrm{R}\) are such that \(\alpha + \beta \in \mathrm{R}\) , then \(\varphi(\alpha + \beta) = \varphi(\alpha) + \varphi(\beta)\) .
\end{enumerate}

Let \(\mathbf{Q}\) be the subgroup of \(\mathbf{V}\) generated by \(\mathbf{R}\) . Then \(\varphi\) extends to a homomorphism from \(\mathbf{Q}\) to \(\Gamma\) .

Let B be a basis of R. Let \(\psi\) be the unique homomorphism from Q to \(\Gamma\) that coincides with \(\varphi\) on B. It suffices to show that \(\psi(\alpha) = \varphi(\alpha)\) when \(\alpha\) is a positive root relative to B. We have \(\alpha = \beta_{1} + \cdots + \beta_{m}\) with \(\beta_{i} \in B\) for all i, and \(\beta_{1} + \cdots + \beta_{h} \in R\) for all h (Prop. 19). We show that \(\psi(\alpha) = \varphi(\alpha)\) by induction on m. This is clear if m = 1. The induction hypothesis gives

\[
\psi \left(\beta_ {1} + \dots + \beta_ {m - 1}\right) = \varphi \left(\beta_ {1} + \dots + \beta_ {m - 1}\right),
\]

and we have \(\psi (\beta_m) = \varphi (\beta_m)\) , hence \(\psi (\alpha) = \varphi (\alpha)\) , which proves the corollary.

For any root \(\alpha = \sum_{\beta \in \mathrm{B}(\mathbb{C})} n_{\beta} \beta\) in \(\mathbb{R}\) , denote by \(\Upsilon(\alpha)\) the set of \(\beta \in \mathrm{B}(\mathbb{C})\) such that \(n_{\beta} \neq 0\) . Moreover, observe that \(\mathrm{B}(\mathbb{C})\) can be identified with the set of vertices of the graph of the Coxeter system formed by \(\mathrm{W}(\mathbb{R})\) and the \(s_{\alpha_i}\) (cf.~Chap. IV, § 1, no. 9 and Chap. V, § 3, no. 2).

COROLLARY 3. a) Let \(\alpha \in \mathbb{R}\) . Then \(\mathrm{Y}(\alpha)\) is a connected subset of \(\mathrm{B}(\mathrm{C})\) (Chap. IV, Appendix).

\begin{enumerate}
\def\labelenumi{\alph{enumi})}
\setcounter{enumi}{1}
\tightlist
\item
  Let Y be a non-empty connected subset of B(C). Then \(\sum_{\beta\in Y}\beta\) belongs to R.
\end{enumerate}

To prove \(a\) ), we can assume that \(\alpha\) is positive. We argue by induction on \(\operatorname{Card}(\mathbf{Y}(\alpha))\) , the assertion being trivial if \(\operatorname{Card}(\mathbf{Y}(\alpha)) = 1\) . By Prop. 19, there exists \(\beta \in \mathrm{B}(\mathrm{C})\) such that \(\alpha - \beta \in \mathrm{R}\) . Let \(p\) be the largest integer \(\geqslant 0\) such that \(\gamma = \alpha - p\beta \in \mathrm{R}\) . Since \(\gamma - \beta \notin \mathrm{R}\) and \(\gamma + p\beta \in \mathrm{R}\) , \((\gamma|\beta) \neq 0\) (Prop. 9); thus \(\beta\) is linked to at least one element of \(\mathbf{Y}(\gamma)\) . But \(\mathbf{Y}(\alpha) = \mathbf{Y}(\gamma) \cup \{\beta\}\) , and \(\mathbf{Y}(\gamma)\) is connected by the induction hypothesis. Thus \(\mathbf{Y}(\alpha)\) is connected, which proves \(a\) ).

Now let Y be a non-empty connected subset of B(C); we show by induction on \(\operatorname{Card}(Y)\) that \(\sum_{\beta\in Y}\beta\) is a root. The case in which \(\operatorname{Card}(Y)\leqslant1\) is trivial. Assume that \(\operatorname{Card}(Y)\geqslant2\) . Since X is a forest (Chap. V, § 4, no. 8, Prop. 8), Y is a tree and has a terminal vertex \(\beta\) (Chap. IV, Appendix). The set \(Y-\{\beta\}\) is connected, and one of its element is linked to \(\beta\) . By the induction hypothesis, \(\alpha=\sum_{\gamma\in Y-\{\beta\}}\gamma\in R\) , and since \((\alpha|\beta)<0\) , it follows that \(\alpha+\beta\in R\) (Th. 1). Q.E.D.

